\documentclass[12pt]{article}
\usepackage[utf8]{inputenc}
\usepackage[T1]{fontenc}
\usepackage{hyperref}
\title{Ubi Language Resources for K-12 Curriculum Development}
\author{Compiled by Claude Sonnet 5 (claude-sonnet-5) for Andr\'as Kornai}
\date{}
\begin{document}
\maketitle

\section{Introduction}
Ubi (also called Oubi) is an Afro-Asiatic language of the East Chadic branch. It is spoken in the Gu\'era region of central Chad, in the Mongo subprefecture near Oubi village. Speaker counts differ between sources. Wikipedia gives about 1,100 (a 1995 figure), while an Ethnologue-based Joshua Project profile gives about 9,000 for the Ubi people. The UNESCO Atlas classes the language as severely endangered. A 2006 SIL survey reportedly found it used at home and in villages, with some shift toward Arabic among younger speakers. The search found very little material that could be used in schools.

\section{Primary Educational Resources}
No primers, textbooks or school materials in Ubi turned up in the searches. This section is empty.

\section{Linguistic and Reference Materials}
\begin{itemize}
\item Wikipedia, Ubi language (overview and classification): \url{https://en.wikipedia.org/wiki/Ubi_language}
\item Max Planck Institute, Chadic numerals database, Ubi page (numeral data): \url{https://lingweb.eva.mpg.de/channumerals/Ubi.htm}
\item Endangered Languages Project, entry for Oubi: \url{https://endangeredlanguages.com/lang/2595?hl=de}
\end{itemize}

\section{Religious and Community Resources}
\begin{itemize}
\item Joshua Project profile, Ubi in Chad (demographic and mission-oriented information): \url{https://joshuaproject.net/people_groups/15679}
\end{itemize}

\section{Digital and Audio Resources}
No audio, video or corpus resources turned up in the searches.

\section{Curriculum Development Context}
Little published material exists for Ubi. A curriculum would have to start from new documentation and orthography work, ideally with community speakers. The reference description says Ubi is taught as a subject in some schools. No source found in the searches documents this. Mother-tongue education in Chad would also need to be checked against Chadian Arabic and French, the languages of wider communication.
\end{document}